\documentclass[conference]{IEEEtran}
\usepackage{amsmath,amssymb}
\usepackage{graphicx}
\usepackage{cite}
\usepackage{float}
\usepackage[hidelinks]{hyperref}
\usepackage{array}

\title{A Hybrid Framework for Jammer Detection Using Statistical Signal Processing and Polarization Features}

\author{\IEEEauthorblockN{Shreyash Puri}
\IEEEauthorblockA{Summer Research Fellow\\
Indian Institute of Technology Madras\\
July 2026}}

\begin{document}

\maketitle

\begin{abstract}
This paper presents a jammer detection framework for wireless communication systems. The problem is formulated as binary hypothesis testing and addressed through three progressively challenging scenarios: (i)~known jammer statistics with Bayesian detection, (ii)~unknown prior probabilities with minimax detection, and (iii)~unknown jammer distribution with residual-energy detection. A polarization-based detector exploiting phase-polarization manifold consistency is developed for strong-jammer scenarios where signal reconstruction fails. The framework is validated through Monte Carlo simulations (achieving ${\sim}22\times$ residual energy separation), Simulink block-level modeling, and real-time hardware implementation using HackRF SDR devices at 915~MHz with a Python-based Costas loop receiver. The results demonstrate that combining statistical signal-domain methods with electromagnetic-domain polarization features provides a robust detection framework across varying jammer strengths.
\end{abstract}

\begin{IEEEkeywords}
Jammer detection, hypothesis testing, likelihood ratio test, residual-energy detection, polarization, software-defined radio, Costas loop.
\end{IEEEkeywords}

%=============================================================================
\section{Introduction}
%=============================================================================

Wireless communication systems are vulnerable to intentional interference, commonly known as jamming. A jammer emits signals within the operating frequency band of a legitimate system to degrade communication quality or disrupt reception. As wireless technologies become critical in civilian, industrial, and defense applications, reliable jammer detection has become an important research problem~\cite{poisel2011,pirayesh2022}.

Conventional detection techniques based on received power or SNR may perform poorly when the jammer is weak or when its behavior resembles normal channel effects~\cite{xu2006}. This work formulates jammer detection as a binary hypothesis testing problem and employs statistical signal processing techniques---likelihood ratio tests (LRT), generalized likelihood ratio tests (GLRT), minimax detection, and matched filtering---to derive optimal decision rules~\cite{kay1998,poor1994}. The approach progresses from the most informed scenario (known jammer statistics and priors) to the least informed (unknown distribution and priors), reflecting the practical reality that adversaries do not reveal their transmission characteristics.

Residual-signal analysis is developed by estimating and removing the legitimate signal from the received waveform. The key insight is that if the signal can be reconstructed, any remaining energy beyond what noise alone explains constitutes evidence of jamming. Furthermore, polarization-based features are utilized, exploiting the vector nature of electromagnetic waves: a jammer is an independent source with its own polarization state that perturbs the expected phase-polarization manifold of the legitimate channel~\cite{born1999,lee1986}.

\subsection{Original Contributions}

The major contributions of this work are:
\begin{enumerate}
\item Development of Bayesian and minimax jammer detectors under progressively relaxed information assumptions, with explicit derivation of optimal thresholds and their dependence on system parameters.
\item Design of a jammer-distribution-free residual-energy detector that requires no knowledge of jammer statistics, achieving ${\sim}22\times$ energy separation between hypotheses in simulation.
\item Investigation of a phase-polarization consistency metric for strong jammer scenarios where signal reconstruction fails, demonstrating complementary performance to residual-based methods.
\item Validation through three levels: Monte Carlo simulation, Simulink block-level modeling, and real-time SDR-based hardware experiments at 915~MHz.
\item Implementation of a complete real-time receiver with Costas loop carrier recovery, signal cancellation, and jammer extraction on HackRF SDR hardware.
\end{enumerate}

%=============================================================================
\section{Detection Framework}
%=============================================================================

The jammer detection problem is decomposed into three progressively challenging cases based on the amount of information available to the receiver. This decomposition follows the classical structure of detection theory~\cite{kay1998,vantrees2001}, where the optimal detector changes fundamentally depending on what the receiver knows about the interference.

\subsection{Bayesian Detection with Known Jammer Statistics}

Consider a binary hypothesis testing problem where the received observation~$y$ is governed by:
\begin{align}
H_0 &: y(t) = s(t) + n(t), \\
H_1 &: y(t) = s(t) + j(t) + n(t),
\end{align}
where $s(t)$ is the legitimate signal, $j(t)$ is the jammer signal, and $n(t)$ denotes additive Gaussian noise.

Assuming the jammer distribution, prior probabilities, and decision costs are known, the optimal Bayesian decision rule is the likelihood ratio test:
\begin{equation}
\Lambda(y) = \frac{f_{Y|H_1}(y)}{f_{Y|H_0}(y)} \underset{H_0}{\overset{H_1}{\gtrless}} \eta,
\end{equation}
where
\begin{equation}
\eta = \frac{P(H_0)\,[C_{10}-C_{00}]}{P(H_1)\,[C_{01}-C_{11}]}.
\end{equation}

For the simplified case of a constant signal $s(t)=S$, constant jammer $j(t)=J$, and $n \sim \mathcal{N}(0,\sigma^2)$, the test reduces to:
\begin{equation}
y \underset{H_0}{\overset{H_1}{\gtrless}} \gamma = S + \frac{J}{2} + \frac{\sigma^2}{J}\ln(\eta).
\end{equation}

The threshold consists of the midpoint $S+J/2$ between hypothesis means, adjusted by $(\sigma^2/J)\ln(\eta)$, which encodes the influence of prior information relative to observation quality.

\subsection{Minimax Detection with Unknown Priors}

When prior probabilities are unknown, the Bayesian threshold cannot be computed. The minimax detector minimizes the worst-case risk by equalizing false alarm and missed detection probabilities:
\begin{equation}
P_{FA} = Q\!\left(\frac{\gamma-S}{\sigma}\right), \quad P_M = Q\!\left(\frac{S+J-\gamma}{\sigma}\right).
\end{equation}

Setting $P_{FA} = P_M$ and exploiting the monotonicity of the Q-function yields:
\begin{equation}
\boxed{\gamma = S + \frac{J}{2}}.
\end{equation}

This threshold is the midpoint between hypothesis means---independent of prior probabilities. Monte Carlo simulation with $10^5$ trials ($A_1=1$, $A_2=2$, $\sigma^2=1$) confirms effective discrimination with the ROC curve lying above the diagonal (Fig.~\ref{fig:roc}).

\begin{figure}[t]
\centering
\includegraphics[width=0.75\columnwidth]{2.2_result.png}
\caption{ROC curve for the LRT statistic $T = j(y-s)$. The minimax operating point corresponds to $P_D = 1 - P_{FA}$.}
\label{fig:roc}
\end{figure}

\subsection{Jammer-Distribution-Free Residual-Based Detection}

When the jammer distribution is unknown, the likelihood ratio cannot be evaluated. The key observation is that while the receiver cannot characterize the jammer, it can accurately model the legitimate signal. If the signal estimate $\hat{s}$ is subtracted from the observation, the residual under $H_0$ contains only noise, while under $H_1$ it contains jammer energy.

The residual sequence is $r_k = y_k - \hat{s}_k$, and the detection statistic is:
\begin{equation}
T_r = \frac{1}{N}\sum_{k=1}^{N} r_k^2.
\end{equation}

Under $H_0$: $E[T_r] = \sigma_n^2$. Under $H_1$: $E[T_r] = P_j + \sigma_n^2$, where $P_j$ is the jammer power. The decision rule is:
\begin{equation}
T_r \underset{H_0}{\overset{H_1}{\gtrless}} \gamma_r.
\end{equation}

Simulation with $N=100$ samples, $s_k=1$, $\sigma=0.3$, and a sinusoidal jammer $j_k = 2\sin(2\pi k/20)$ yields residual energies:
\begin{equation}
E_{H_0} = 9.44, \quad E_{H_1} = 209.76, \quad \frac{E_{H_1}}{E_{H_0}} \approx 22.2\times.
\end{equation}

This large separation validates the detector's effectiveness without any knowledge of jammer statistics (Fig.~\ref{fig:residual}).

\begin{figure}[t]
\centering
\includegraphics[width=0.48\columnwidth]{receivednojampng}
\includegraphics[width=0.48\columnwidth]{resduenojampng}\\[3pt]
\includegraphics[width=0.48\columnwidth]{receivedjam}
\includegraphics[width=0.48\columnwidth]{resduejam}
\caption{Residual-based detection. Top row: received signal and residual under $H_0$. Bottom row: received signal and residual under $H_1$.}
\label{fig:residual}
\end{figure}

\subsubsection{Limitations}
When $P_j \gtrsim P_s$, signal estimation degrades because the jammer corrupts demodulation. The residual then contains both jammer energy and estimation errors, reducing reliability. This motivates polarization-based methods.

%=============================================================================
\section{Polarization-Based Jammer Detection}
%=============================================================================

\subsection{Motivation}

Polarization-based detection exploits the vector nature of electromagnetic waves. A legitimate transmitter has a known polarization state, multipath produces predictable polarization changes, and a jammer introduces polarization inconsistencies that cannot be explained by the channel model alone~\cite{born1999}.

\subsection{Multipath Polarization Model}

For a vertically polarized transmission with a direct path and a reflected path (two-ray model), the field components are:
\begin{align}
E_H &= \alpha \sin(\beta)\,e^{j\Delta\phi}, \\
E_V &= 1 + \alpha \cos(\beta)\,e^{j\Delta\phi},
\end{align}
where $\alpha$ is the multipath attenuation, $\beta$ is the polarization rotation, and $\Delta\phi$ is the path phase difference. The expected polarization ratio is:
\begin{equation}
Q_{\mathrm{exp}} = \frac{|E_V|^2}{|E_H|^2 + \varepsilon}.
\end{equation}

For $\alpha=2$ and $\beta=5^\circ$, the maximum polarization deviation is $\Delta\psi_{\max} = 9.96^\circ$ at $\Delta\phi \approx 180^\circ$ (Fig.~\ref{fig:pol1}).

\begin{figure}[t]
\centering
\includegraphics[width=0.8\columnwidth]{pol1.png}
\caption{Phase-polarization manifold showing deterministic relationship between multipath phase shift and polarization deviation.}
\label{fig:pol1}
\end{figure}

\subsection{Phase-Polarization Consistency Metric}

A jammer with amplitude $A_j$, polarization angle $\theta_j$, and phase $\phi_j$ perturbs the observed manifold. The consistency metric is:
\begin{equation}
u_k = 10\left|\log_{10}\!\left(\frac{Q_{\mathrm{obs}}(\Delta\phi_k)}{Q_{\mathrm{exp}}(\Delta\phi_k)}\right)\right|,
\end{equation}
with the detection statistic $\bar{u} = (1/N)\sum_{k=1}^{N} u_k$. The decision rule is:
\begin{equation}
\bar{u} \underset{H_0}{\overset{H_1}{\gtrless}} \gamma.
\end{equation}

\subsection{Simulation Results}

Monte Carlo simulations (5000 trials) under Rayleigh fading and AWGN yield:
\begin{itemize}
\item Clear separation between $H_0$ and $H_1$ distributions (Fig.~\ref{fig:pol2}).
\item ROC above the diagonal with $P_D = 0.4354$ at $P_{FA} = 0.1$ (Fig.~\ref{fig:pol3}).
\item Detection probability increases with JSR: $P_D \approx 0.09$ at $-20$~dB, $P_D \approx 0.63$ at $+10$~dB (Fig.~\ref{fig:pol4}).
\end{itemize}

\begin{figure}[t]
\centering
\includegraphics[width=0.8\columnwidth]{pol2.png}
\caption{PDF of the consistency metric under $H_0$ and $H_1$.}
\label{fig:pol2}
\end{figure}

\begin{figure}[t]
\centering
\includegraphics[width=0.8\columnwidth]{pol3.png}
\caption{ROC curve of the polarization-based detector.}
\label{fig:pol3}
\end{figure}

\begin{figure}[t]
\centering
\includegraphics[width=0.8\columnwidth]{pol4.png}
\caption{Detection probability versus jammer-to-signal ratio.}
\label{fig:pol4}
\end{figure}

\subsection{Complementary Performance}

Residual-based detection works well for weak jammers ($P_j \ll P_s$), while polarization-based detection improves as jammer power grows ($P_j \gg P_s$). This complementarity motivates a hybrid framework that adaptively combines both mechanisms.

%=============================================================================
\section{Simulink Implementation}
%=============================================================================

A Simulink block-level model validates the residual-based detector in a streaming environment. The processing chain implements: BPSK generation $\rightarrow$ AWGN channel $\rightarrow$ jammer injection $\rightarrow$ matched filtering $\rightarrow$ signal subtraction $\rightarrow$ energy computation $\rightarrow$ threshold comparison (Fig.~\ref{fig:simulink}).

\begin{figure}[t]
\centering
\includegraphics[width=\columnwidth]{simulink.png}
\caption{Simulink block diagram of the residual-based jammer detection system.}
\label{fig:simulink}
\end{figure}

Key observations from the Simulink model:
\begin{enumerate}
\item The matched filter provides accurate signal estimation under combined noise and jammer interference.
\item Signal cancellation isolates the jammer component in the residual.
\item The squared residual energy is consistently elevated when a jammer is present, confirming $E[r^2] = P_j + \sigma_n^2$.
\item An oscillatory pattern in the residual provides a structural detection criterion complementing energy thresholding.
\end{enumerate}

A two-stage detection logic is implemented: Stage~1 examines the residual for oscillatory structure (detecting structured jammers), and Stage~2 monitors residual power against the noise baseline (detecting noise jammers).

%=============================================================================
\section{Hardware Implementation}
%=============================================================================

\subsection{Experimental Setup}

The hardware testbed uses three HackRF One SDR devices operating at $f_c = 915$~MHz with $F_s = 2$~MHz sample rate. The transmitter sends a pre-generated BPSK IQ file, a second HackRF transmits the jammer waveform, and a third HackRF receives via a Python-based real-time receiver.

\subsection{Receiver Architecture}

The receiver implements:
\begin{enumerate}
\item \textbf{IQ acquisition}: 8-bit samples streamed at 2~MHz via \texttt{hackrf\_transfer} in pipe mode. Frame size: 500,000 samples (0.25~s).
\item \textbf{Frequency acquisition}: FFT-based estimation with parabolic interpolation for sub-bin accuracy.
\item \textbf{Costas loop carrier recovery}: Second-order loop ($B_n = 1000$~Hz, $\zeta = 0.707$) with Numba JIT compilation for real-time performance. The loop updates are:
\begin{align}
\omega_{n+1} &= \omega_n + \beta \cdot e(n), \\
\phi_{n+1} &= \phi_n + \omega_{n+1} + \alpha \cdot e(n),
\end{align}
where the BPSK error signal is $e(n) = \mathrm{sgn}(I_{\mathrm{out}})\cdot Q_{\mathrm{out}}$.
\item \textbf{Signal reconstruction}: Hard decision followed by re-modulation using the tracked carrier phase.
\item \textbf{Jammer extraction}: Signal subtraction and 4th-order Butterworth bandpass filtering (500~Hz to 100~kHz).
\item \textbf{Power measurement}: $P_j = 10\log_{10}(\mathrm{mean}(r_I^2 + r_Q^2))$~dBFS.
\end{enumerate}

\subsection{Experimental Results}

Fig.~\ref{fig:results} shows the receiver output with a 1~kHz sinusoidal jammer at 30\% signal power. The Costas loop maintains lock on the BPSK carrier, the extracted jammer waveform shows a clear sinusoidal pattern after signal cancellation, and the residual power is elevated ${\sim}30$~dB above the noise floor---demonstrating successful detection under real-world conditions including frequency offset, phase noise, and 8-bit quantization.

\begin{figure}[t]
\centering
\includegraphics[width=\columnwidth]{aftertran.png}
\caption{Hardware results with sinusoidal jammer active. Top: stable BPSK demodulation. Middle: extracted 1~kHz jammer. Bottom: elevated jammer power.}
\label{fig:results}
\end{figure}

%=============================================================================
\section{Conclusion}
%=============================================================================

This work developed a comprehensive jammer detection framework progressing from Bayesian detection through minimax detection to jammer-distribution-free residual-energy detection. A polarization-based detector addresses the strong-jammer regime where signal estimation degrades. The framework was validated at three levels: Monte Carlo simulation, Simulink modeling, and real-time SDR hardware implementation at 915~MHz.

The key results are: (i)~the residual-energy detector achieves ${\sim}22\times$ energy separation between hypotheses, (ii)~the polarization detector provides complementary performance at high JSR, and (iii)~the hardware receiver successfully extracts a 1~kHz sinusoidal jammer after Costas loop demodulation and signal cancellation.

\begin{table}[t]
\centering
\caption{Comparison of Detection Approaches}
\label{tab:comparison}
\begin{tabular}{|l|c|c|c|}
\hline
\textbf{Method} & \textbf{Needs jammer} & \textbf{Works at} & \textbf{Hardware} \\
 & \textbf{statistics?} & \textbf{high JSR?} & \textbf{feasible?} \\
\hline
Bayesian/Minimax & Yes & Medium & Yes \\
\hline
Residual-energy & No & Weak & Yes \\
\hline
Polarization & No & Strong & Harder \\
\hline
\end{tabular}
\end{table}

Future work includes hardware validation of the polarization detector using dual-polarized antennas, adaptive threshold selection for non-stationary environments, and evaluation under frequency-hopping and spread-spectrum jamming scenarios.

%=============================================================================
% REFERENCES
%=============================================================================

\begin{thebibliography}{21}

\bibitem{kay1998}
S.~M.~Kay, \textit{Fundamentals of Statistical Signal Processing, Volume~II: Detection Theory}, Prentice Hall, 1998.

\bibitem{vantrees2001}
H.~L.~Van Trees, \textit{Detection, Estimation, and Modulation Theory, Part~I}, John Wiley \& Sons, 2001.

\bibitem{poor1994}
H.~V.~Poor, \textit{An Introduction to Signal Detection and Estimation}, 2nd~ed., Springer, 1994.

\bibitem{poisel2011}
R.~A.~Poisel, \textit{Modern Communications Jamming: Principles and Techniques}, 2nd~ed., Artech House, 2011.

\bibitem{proakis2008}
J.~G.~Proakis and M.~Salehi, \textit{Digital Communications}, 5th~ed., McGraw-Hill, 2008.

\bibitem{goldsmith2005}
A.~Goldsmith, \textit{Wireless Communications}, Cambridge University Press, 2005.

\bibitem{born1999}
M.~Born and E.~Wolf, \textit{Principles of Optics}, 7th~ed., Cambridge University Press, 1999.

\bibitem{scharf1991}
L.~L.~Scharf, \textit{Statistical Signal Processing: Detection, Estimation, and Time Series Analysis}, Addison-Wesley, 1991.

\bibitem{haykin2001}
S.~Haykin, \textit{Communication Systems}, 4th~ed., John Wiley \& Sons, 2001.

\bibitem{simon2005}
M.~K.~Simon, S.~M.~Hinedi, and W.~C.~Lindsey, \textit{Digital Communication Techniques: Signal Design and Detection}, Prentice Hall, 1995.

\bibitem{lathi2009}
B.~P.~Lathi and Z.~Ding, \textit{Modern Digital and Analog Communication Systems}, 4th~ed., Oxford University Press, 2009.

\bibitem{xu2006}
W.~Xu, W.~Trappe, Y.~Zhang, and T.~Wood, ``The feasibility of launching and detecting jamming attacks in wireless networks,'' in \textit{Proc. ACM MobiHoc}, 2005, pp.~46--57.

\bibitem{mpitziopoulos2009}
A.~Mpitziopoulos, D.~Gavalas, C.~Konstantopoulos, and G.~Pantziou, ``A survey on jamming attacks and countermeasures in WSNs,'' \textit{IEEE Commun. Surveys Tuts.}, vol.~11, no.~4, pp.~42--56, 2009.

\bibitem{pirayesh2022}
H.~Pirayesh and H.~Zeng, ``Jamming attacks and anti-jamming strategies in wireless networks: A comprehensive survey,'' \textit{IEEE Commun. Surveys Tuts.}, vol.~24, no.~2, pp.~767--809, 2022.

\bibitem{lin2005}
G.~Lin and G.~Noubir, ``On link layer denial of service in data wireless LANs,'' \textit{Wireless Commun. Mobile Comput.}, vol.~5, no.~3, pp.~273--284, 2005.

\bibitem{gardner1988}
W.~A.~Gardner, ``Signal interception: A unifying theoretical framework for feature detection,'' \textit{IEEE Trans. Commun.}, vol.~36, no.~8, pp.~897--906, 1988.

\bibitem{rice2004}
M.~Rice, \textit{Digital Communications: A Discrete-Time Approach}, Prentice Hall, 2004.

\bibitem{lindsey1972}
W.~C.~Lindsey and M.~K.~Simon, \textit{Telecommunication Systems Engineering}, Prentice Hall, 1973.

\bibitem{skolnik2008}
M.~I.~Skolnik, \textit{Radar Handbook}, 3rd~ed., McGraw-Hill, 2008.

\bibitem{lee1986}
J.~S.~Lee and E.~Pottier, \textit{Polarimetric Radar Imaging: From Basics to Applications}, CRC Press, 2009.

\bibitem{ossmann2014}
M.~Ossmann, ``HackRF One: An open source SDR platform,'' Great Scott Gadgets, 2014.

\end{thebibliography}

\end{document}
